\section{The compact bound from a uniformly generated family}
\label{c:section}

The compact instance uses the component theorems of Supplement~C
of~\cite{BoundsSupplement}. We state their energy bounds and complete
domains; the common LP then performs the global accounting.

Fix the original farthest vertex and the original vector $x$ used in
Theorem~\ref{thm:energy-normalization}. If $C$ is an irreducible spherical
component of its vertex diagram, put
\[
 e_C=\norm{\operatorname{proj}_{\langle C\rangle}x}^2.
\]
The component subspaces are mutually orthogonal and exhaust the spatial
space. Consequently
\begin{equation}\label{c:budget}
 e_C\ge0,\qquad \sum_C e_C=\norm{x}^2<1.
\end{equation}
For disjoint collections $\mathcal S$ of components we write
$e_{\mathcal S}=\sum_{C\in\mathcal S}e_C$. Throughout this section, a
\emph{spectator} is another component at this same vertex, with its actual
opposite walls retained. It is not an independently chosen diagram.

Vinberg's compact dimension theorem~\cite{Vinberg1984} excludes dimensions
at least $30$. The LP covers each actual dimension $18\le n\le29$;
no passage to a lower-dimensional face is used.

\subsection{Small components and the actual exceptional profiles}
\label{c:small}

We use the usual notation for irreducible spherical Coxeter diagrams.
The rank-two family is $I_2(m)$, where $m\ge3$ is an integer. We identify
$A_2$ with $I_2(3)$ and do not count aliases such as $B_2$ or $D_3$ a
second time. An \emph{ordinary} component is a component that does not
carry one of the actual exceptional wall profiles described below.

A profile specifies an opposite wall of the given polytope and its
contacts with the original walls. An old bond labelled seven alone
does not specify $O_7$. All labels below belong to pairs of actual
original walls. Auxiliary projected vectors are not assigned those labels.

\begin{proposition}[Ordinary component estimates, retained input]
\label{c:ordinary-input}
For the same original vector $x$, the ordinary components in
Table~\ref{c:ordinary-table} have energy strictly greater than the listed
number. Every estimate permits arbitrary finite spectators.
\end{proposition}

\begin{table}[htbp]
\centering
\caption{The ordinary component lower bounds. Each decimal in this paper
is an exact rational number.}\label{c:ordinary-table}
\begin{tabular}{ll@{\qquad}ll}
\toprule
Component&Lower bound&Component&Lower bound\\\midrule
$A_1$&$1/10$&$A_2$&$319/2000$\\
$I_2(m),\ m\ge4$&$1/7$&$A_3$&$1/4$\\
$B_3,H_3$&$2/7$&$A_4$&$2/7$\\
$B_4,D_4,F_4$&$1/3$&$H_4$&$11/36$\\
$A_5$&$3/8$&$D_5$&$5/14$\\
$A_6,B_5,B_6,D_6$&$3/7$&$B_7,D_7$&$1/2$\\
$B_8,D_8$&$4/7$&$E_6$&$9/20$\\
$E_7$&$5/9$&$E_8$&$2/3$\\\bottomrule
\end{tabular}
\end{table}

These are the ordinary-component conclusions of Supplement~C, including
the $B_3,H_3$ improvement. The dihedral row quantifies over every integer
$m\ge4$, rather than a finite sample of labels.

For clarity, we describe the exceptional profiles before using their
energy debits. Let $p$ be the original unit timelike vertex vector, and
let $u$ be the unit normal of an original $A_1$ component. An isolated
short rank-one opposite has the form
\begin{equation}\label{c:I-normal}
 c_I=\sqrt h\,p-\sqrt{1+h}\,u,
 \qquad 0<h<\frac12.
\end{equation}
Its spatial support is this one $A_1$. This is profile $I$.
Profile $J$ is the distinct short-wall alternative supported on two old
$A_1$ factors, with a dotted contact on one and an actual contact labelled
three on the other. Its squared time lies in $(1/3,3/8)$. The two factors
belong to one profile and must be accounted for jointly.

There are also two rank-two profiles:
\[
 N_7=(\text{old }A_2,\ \text{actual new isolated contact }7),
 \qquad
 O_7=(\text{old }I_2(7),\ \text{actual new isolated contact }3).
\]
Here isolation means that the new wall has no spatial contact with
another old component. These specifications include the new wall, not
merely the type of its old support.

\begin{proposition}[Exceptional profiles and compatibility, retained input]
\label{c:defect-input}
At most one of $I,J,N_7,O_7$ occurs, and the alternatives do not coexist.
Their owned ranks and energy bounds are
\[
\begin{array}{c|c|c}
\text{profile}&\text{old support}&\text{energy lower bound}\\\hline
I&A_1&0\\
J&A_1+A_1&17/50\\
N_7&A_2&9/100\\
O_7&I_2(7)&9/100.
\end{array}
\]
The positive entries are strict in the geometric theorem. No owned
component receives an ordinary bound in addition to its profile bound.
\end{proposition}

Compatibility is part of the actual-wall theorem, including overlapping
short-wall alternatives; it is not inferred from the spherical types
of the supports.

Two further input families will be used. If $T_j$ is an ordinary
irreducible of rank $j$, then
\begin{align}
 e_{I+A_1}&>1/5,&
 e_{I+2A_1}&>3/10,\label{c:I-joint-one}\\
 e_{I+T_2}&>21/100,&
 e_{I+T_3}&>2/7.\label{c:I-joint-two}
\end{align}
Each left side is one joint support. Overlapping uses of $I$ therefore
share its single capacity in Theorem~\ref{thm:energy-certificate}.

Beside an actual $N_7$ or $O_7$, every other ordinary rank-two component
has energy greater than $1/5$. An additional ordinary $A_1$ has energy
greater than $1/5$ beside $N_7$, and greater than $2/15$ beside $O_7$.
These propagated estimates concern the additional component alone.
Their domains and the joint bounds above are those of the
ordinary-component appendix of Supplement~C.

\subsection{The larger component families}
\label{c:large}

The $A$-component theorems give
\begin{equation}\label{c:A-thresholds}
 e_{A_r}>\frac12\quad(r\ge7),\qquad
 e_{A_r}>\frac47\quad(r\ge8),\qquad
 e_{A_r}>\frac57\quad(r\ge13).
\end{equation}
All three estimates allow arbitrary finite spectators. The exact-rank
improvements retained by the LP are listed in
Table~\ref{c:A-table}. These are bounds on the original projection,
even when their local proofs use points or neighboring vertices.

\begin{table}[htbp]
\centering
\caption{Additional spectator-uniform $A$-component bounds.}
\label{c:A-table}
\begin{tabular}{clp{.57\linewidth}}
\toprule
Rank&Lower bound&Local result in Supplement~C\\\midrule
12&$13/20$&The two-end $A_{12}$ estimate\\
15&$413/500$&The $A_{15}$ reference-vector estimate\\
16&$43/50$&The $N=17$ reference-vector estimate\\
17&$6/7$&The $A$-chain central-window argument\\
18&$91/100$&The $N=19$ chain window and reference-vector estimate\\
19&$563/625$&The $A_{19}$ projection-energy lemma\\
20&$.92274$&The $A_{20}$ projection-energy lemma\\\bottomrule
\end{tabular}
\end{table}

For actual $B_r,D_r$ components with $r\ge9$, the common bound is
\begin{equation}\label{c:BD-nine}
                         e_{B_r},e_{D_r}>\frac9{14}.
\end{equation}
The next family uses an initial segment of the original component. Its
rank conditions are part of the statement and will be retained literally.

\begin{proposition}[The generic prefix estimate, retained input]
\label{c:prefix-input}
Let $6\le M\le16$. Suppose that the actual component is $B_r$ with
$r\ge M+5$, or $D_r$ with $r\ge M+6$. Put
\[
 f_i=\sqrt i-\frac{61}{100\sqrt2},\quad
 E_M=\frac{f_5^2}{5}+
        \sum_{i=5}^{M-1}(\sqrt{i+1}-\sqrt i)^2,
 \quad H_M=\left\lfloor\frac{M-5}{2}\right\rfloor.
\]
On the same original vector, the first $M+1$ coordinates of that
component have energy greater than
\begin{equation}\label{c:prefix-formula}
                  E_M-H_M/3200+3969/800000.
\end{equation}
The estimate permits arbitrary finite spectators.
\end{proposition}

The first two terms are the generic prefix estimate. The last term is
the additional initial-coordinate gain proved under exactly the same
hypotheses. It replaces the weaker bound on that prefix; it is not an
independent energy contribution. The rational lower bounds used in the
current program are $P_M+3969/800000$, with $P_M$ as follows:
\begin{center}
\begin{tabular}{c|ccccc}
$M$&7&8&9&10&11\\\hline
$P_M$&$.7349975$&$.7683275$&$.797425$&$.823735$&$.8472225$\\[2pt]
\end{tabular}
\par\smallskip
\begin{tabular}{c|ccccc}
$M$&12&13&14&15&16\\\hline
$P_M$&$.8689525$&$.88864$&$.90715$&$.9240775$&$.9401975$
\end{tabular}
\end{center}
The inequalities $P_M<E_M-H_M/3200$ are the rational evaluations in
Supplement~C. The current finite LP uses $7\le M\le16$ from this family.
For every applicable $M$ it may keep both the weaker and the stronger
row. On identical support and domain the weaker one is redundant, but
retaining it does not alter the validity of the stronger row.

There are three further local continuations. They are stated separately
because their proofs extend a short end where the generic theorem alone
does not apply:
\begin{enumerate}[label=\textup{(\roman*)}]
\item On the entire forest $B_{16}+T_2$ or $D_{16}+T_2$ in actual
 dimension $18$, where $T_2$ is ordinary, the large component has energy
 greater than $.8689525$.
\item On the entire forest $B_{17}+A_1$ in dimension $18$, with ordinary
 nonshort $A_1$, the $B_{17}$ component has energy greater than $.90715$.
\item On the entire forest $D_{19}+A_1$ in dimension $20$, with ordinary
 nonshort $A_1$, the actual $D_4$ continuation and initial gain give
 $e_{D_{19}}>.91211125$.
\end{enumerate}
The first two are the exact-rank-eighteen layer estimates. The third
uses the proper $D_{19}$ continuation before the rank-twenty accounting.
We import that local continuation and its energy calculation, not the
conclusion that dimension twenty is impossible. A bound proved for an
ordinary $T_2$ is not used for $N_7$ or $O_7$, and an ordinary $A_1$
bound is not transferred to $I$.

\subsection{A complete domain construction}
\label{c:cover}

We call $A_r$, $r\ge7$, and $B_r,D_r$, $r\ge9$,
\emph{dominant}. This is a definition for the present accounting, not
a new class of Coxeter diagrams. Every dominant costs more than one
half by \eqref{c:A-thresholds} and \eqref{c:BD-nine}. Every other
ordinary irreducible costs at least its rank divided by fourteen, as
one checks directly from Table~\ref{c:ordinary-table}.

\begin{lemma}[The common spectator estimate]\label{c:spectator}
Let a collection of nondominant components have total rank $S$. Include
every exceptional profile owned by that collection. Its energy satisfies
\begin{equation}\label{c:spectator-rate}
                          e_{\mathcal S}\ge\frac{S-1}{14}.
\end{equation}
For no profile or profile $J$, the stronger bound $S/14$ is available.
For profile $I$ the possible loss is $1/14$, and for $N_7$ or $O_7$
it is $37/700$.
\end{lemma}
\begin{proof}
Sum the ordinary rank-over-fourteen estimates on the unmarked components.
Profile $I$ replaces the ordinary rank-one contribution by a nonnegative
energy, losing at most $1/14$. Profile $J$ owns rank two and its bound
$17/50$ is greater than $2/14$, so it creates no loss. Either typed
rank-two profile has cost at least $9/100$, whose deficit from $2/14$ is
$37/700<1/14$. Proposition~\ref{c:defect-input} permits at most one of
these losses. This proves every assertion.
\end{proof}

For a fixed dimension $18\le n\le29$, the following construction produces
a complete necessary family of energy domains.

\begin{enumerate}[label=\textup{\arabic*.},leftmargin=*]
\item If there is no dominant, use one aggregate energy variable for all
 components and the row $e\ge(n-1)/14$.
\item If there are at least two dominants, select two distinct actual
 dominant components. Give each its own energy variable and the bound
 $1/2$. A third aggregate variable holds every remaining component.
 The selected components' individual ranks need not be fixed for this
 domain: only their disjointness and their common lower bound are used.
\item If there is exactly one dominant of type $T$ and rank $r$, let
 $S=n-r$ be the total spectator rank. Retain an explicit local
 component exclusion whenever one applies. Otherwise, for $S\ge8$ use
 the two rows $e_T\ge1/2$ and $e_{\mathcal S}\ge(S-1)/14$.
\item For $0<S<8$, let $d(T,r)$ be the largest applicable lower bound
 among the spectator-uniform component estimates in
 Section~\ref{c:large}. Only unconditional component estimates are used
 to choose $d(T,r)$; an estimate requiring one particular spectator
 does not qualify. If
 \begin{equation}\label{c:aggregate-test}
                  d(T,r)+(S-1)/14\ge1,
 \end{equation}
 retain the complete domain with the two aggregate rows giving these
 two bounds.
\item In every remaining case, retain every nondominant spherical
 spectator forest of rank $S$ and every allowed assignment of one
 exceptional profile. On that detailed domain include all applicable
 ordinary, component, joint and propagated bounds, together with the
 precisely matched local exclusions described in Section~\ref{c:local-exits}.
 For $S=0$ this step simply gives the whole dominant component.
\end{enumerate}

Coarse domains remain in the LP. Their displayed rows have a certificate
with weights one, including the equality case $S=8$. The test
\eqref{c:aggregate-test} decides whether to retain detailed spectator
types; it does not enlarge any local theorem's domain.

\begin{proposition}[Completeness of the compact family]
\label{c:complete-cover}
Every hypothetical compact Coxeter polytope of actual dimension
$18\le n\le29$, with its chosen original farthest vertex, gives either
one of the explicitly retained local contradictions or an energy point
in a domain of the above construction. Every coordinate is an energy
of the original vector, and every row applies on the entire domain to
which it is attached.
\end{proposition}
\begin{proof}
Spherical classification gives zero, at least two, or exactly one
dominant component. The first two alternatives give steps one and two.
In the third, its type and rank fix $S$; Lemma~\ref{c:spectator} and the
applicable component bound justify either coarse description.
Otherwise step five includes the actual spectator multiset and its
unique allowed profile. Equal-type components may be relabelled with
their actual walls to put a marked component first. This preserves all
other parameters: an ordinary dihedral component retains its full
integer range, while an $O_7$ mark fixes only its owned component to
label seven. Every row and local exclusion is then used with its stated
component, profile and spectator hypotheses. Orthogonality gives
\eqref{c:budget} on the original vector.
\end{proof}

To generate a detailed remainder, use this ordered list of nondominant
types through spectator rank seven:
\[
\begin{array}{c|l}
\text{rank}&\text{types}\\\hline
1&A_1\\
2&A_2=I_2(3),\ I_2(m)\ (m\ge4)\\
3&A_3,B_3,H_3\\
4&A_4,B_4,D_4,F_4,H_4\\
5&A_5,B_5,D_5\\
6&A_6,B_6,D_6,E_6\\
7&B_7,D_7,E_7.
\end{array}
\]
Choose successive types in nondecreasing order, subtracting their ranks
until the remainder is zero. Then mark none, one $A_1$ as $I$, two
$A_1$ factors as $J$, one $A_2$ as $N_7$, or one dihedral factor as
$O_7$. Nondecreasing order and exchangeability remove repetitions,
while all dihedral parameters remain quantified. The construction
produces the $829$ domains checked below; its completeness follows
from Proposition~\ref{c:complete-cover}, independently of that count.

\subsection{The retained local contradiction domains}
\label{c:local-exits}

Table~\ref{c:local-exit-table} gives the local contradiction inputs of
Supplement~C. Their normalized mass is zero; no component energy floor
is inferred from them. ``Entire'' specifies the complete original
vertex forest. Apart from C1, only the displayed entire-forest domains
are used, with every Coxeter label and separated-wall parameter in the
cited theorem retained.

\begin{longtable}{>{\raggedright\arraybackslash}p{.08\linewidth}
 >{\raggedright\arraybackslash}p{.43\linewidth}
 >{\raggedright\arraybackslash}p{.39\linewidth}}
\caption{Compact local contradiction inputs. The descriptions refer to
the complete local proofs in Supplement~C.}\label{c:local-exit-table}\\
\toprule
Code&Matched domain&Local proof retained\\\midrule
\endfirsthead
\toprule Code&Matched domain&Local proof retained\\\midrule\endhead
C1&Any $A_r$, $21\le r\le29$; $B_r$, $23\le r\le29$;
 or $D_r$, $24\le r\le29$, with arbitrary finite spectators&
 Large-component cap theorem\\
C2&Entire $B_r+I$, $r\ge20$, or $D_r+I$, $r\ge21$&
 The uniform proper classical-component argument beside $I$\\
C3&Entire $A_{16}+N_7$ or $A_{16}+O_7$&
 The two typed endpoint families in actual dimension $18$\\
C4&Entire $A_{17}+A_1$ or $A_{17}+I$, with ordinary $A_1$&
 The complete rank-one-remainder argument in dimension $18$\\
C5&Entire $B_{16}+N_7/O_7$ or $D_{16}+N_7/O_7$&
 The typed rank-sixteen arguments, retaining both signed $D$ forks\\
C6&Entire $B_{17}+I$&The short rank-one-remainder head argument\\
C7&Entire $D_{17}+A_1$ or $D_{17}+I$, with ordinary $A_1$&
 The separate ordinary and isolated-short rank-one arguments\\
C8&Entire $A_{18}$&The two-end argument using the gap between
 allowed actual wall products\\
C9&Entire $B_{18}$&The full terminal-wall partition, including common
 and distinct opposite walls\\
C10&Entire $D_{18}$&The complete original end-wall and fork argument,
 including all finite, parallel and dotted contacts at the distinguished end\\
C11&Entire $A_{17}+N_7/O_7$ or $A_{18}+I$&
 The separate exact-rank-nineteen local arguments\\
C12&Entire $D_{17}+T_2$, $B_{17}+N_7/O_7$, or
 $B_{18}/D_{18}+A_1$, with ordinary $T_2,A_1$&
 The proper classical rank-nineteen continuation\\
C13&Entire $D_{17}+N_7/O_7$, $B_{18}+I$, or $D_{18}+I$&
 The respective typed and short-remainder arguments at rank nineteen\\
C14&Entire $A_{19}$, $B_{19}$ or $D_{19}$&
 Their separate complete endpoint, head and fork proofs\\
C15&Entire $A_{19}+I$&The $A_{19}$ isolated-short-remainder argument\\
C16&Entire $B_{18}/D_{18}+N_7/O_7$ or $B_{19}/D_{19}+I$&
 The proper classical rank-twenty continuation\\
C17&Entire $A_{20}$, $B_{20}$ or $D_{20}$&
 Their separate whole-component arguments at actual rank twenty\\
C18&Entire $A_{19}+N_7/O_7$ or $A_{20}+I$&
 The typed and isolated-short $A$-component arguments\\
C19&Entire $D_{19}+N_7/O_7$ or $D_{20}+I$&
 The proper $D$-component continuation at rank twenty-one\\
C20&Entire $B_{21}$ or $D_{21}$&
 The separate full-component arguments at actual rank twenty-one\\
C21&Entire $B_{22}$, $D_{22}$ or $D_{23}$&
 The retained full-component terminal argument\\\bottomrule
\end{longtable}

Each entry denotes the complete cited local theorem: C10, for example,
includes all whole-$D_{18}$ end-wall and fork alternatives. A comparison
with another vertex within such a proof contributes no further energy
to the original budget \eqref{c:budget}.

\subsection{The resulting compact certificate}
\label{c:certificate}

The construction gives $829$ branch domains. Of these, $151$ carry one
of the retained local contradictions. On each remaining branch the LP
has one nonnegative multiplier for every applicable energy row. The
capacity at each original projection is one. The exact rational
certificate satisfies the hypotheses of
Theorem~\ref{thm:energy-certificate} simultaneously in every branch.
In particular, overlapping rows involving $I$ share the same $I$ column;
the program cannot buy two copies of its energy.

\begin{theorem}[Compact recovery]\label{c:recovery}
The local geometric inputs stated in Sections~\ref{c:small},
\ref{c:large} and~\ref{c:local-exits}, together with spherical
classification, give an infeasible normalized energy LP for the complete
union of actual dimensions $18,\ldots,29$. With Vinberg's external
compact ceiling $29$, they recover the dimension bound $n\le17$.
\end{theorem}
\begin{proof}
Proposition~\ref{c:complete-cover} assigns every hypothetical polytope
to a complete branch. A branch with a retained local contradiction is
impossible. On every other branch, the exact capacity certificate gives
total original energy at least one by
Theorem~\ref{thm:energy-certificate}, contrary to \eqref{c:budget}.
The exact check covers all $3157$ dual variables and $2224$ rows;
the reproduction section identifies the certificate and model records.
The external theorem excludes dimensions $n\ge30$.
\end{proof}

Thus the LP performs the final rank and energy accounting from valid
local bounds. Its domain construction begins with spherical
classification, without assuming the historical eleven-forest
rank-eighteen list.

\section{The compact arithmetic bound with the same energy program}
\label{a:section}

For a compact arithmetic polytope, its reflection group is an arithmetic
lattice. Supplement~A of~\cite{BoundsSupplement} supplies the further
local restrictions below, with the actual defining field and original
root lines retained. The energy budget remains \eqref{c:budget}.
The compact catalogue already covers dimensions $18,\ldots,29$;
we append the arithmetic domains at actual ranks $16$ and $17$ to
that same catalogue.

\subsection{The retained arithmetic inputs}
\label{a:inputs}

We first recall what the field conditions refer to. In the quadratic-form
model supplied by the arithmetic companion, $K$ is a totally real number
field, $q$ is a quadratic form over $K$, and $B$ is its associated symmetric
bilinear form. One real embedding $\sigma_0:K\hookrightarrow\R$ is
distinguished: at this embedding $q$ has Lorentzian signature and defines
the given hyperbolic space. At every other real embedding $\sigma$, the
conjugate form $q^\sigma$ is positive definite. This is a statement about
the form, not about independence of an arbitrary collection of walls.

The roots $\alpha_i$ span the original wall lines over this same field.
Their Cartan entries are
\begin{equation}\label{a:cartan}
 C_{ij}=\frac{2B(\alpha_i,\alpha_j)}{B(\alpha_i,\alpha_i)}.
\end{equation}
The arithmetic integrality condition used in the local proofs requires
the cyclic products
\begin{equation}\label{a:cycles}
             C_{i_1i_2}C_{i_2i_3}\cdots C_{i_\ell i_1}
\end{equation}
to lie in $\mathcal O_K$, the ring of algebraic integers in $K$.
This formulation is insensitive to changing a root representative on
its line: under $\alpha_i\mapsto t_i\alpha_i$, with $t_i\in K^\times$,
the entry $C_{ij}$ is
multiplied by $t_j/t_i$, and those factors cancel around the cycle.
The companion proves the field reductions on their stated domains.

At a nonphysical embedding $\sigma\ne\sigma_0$, a Gram matrix of
linearly independent conjugate roots is positive definite. A Gram matrix
of all wall roots can have linear dependencies and is only positive
semidefinite. In particular, one must not demand positive definiteness
of a full redundant wall matrix. Similarly, normalizing roots to unit
length can introduce square roots and enlarge the field. The field of
the resulting real cosine entries is not a substitute for the original
common field, root lines and norm classes. The local field-square tests
retain that information when they restrict the possible wall parameters.
Here a root's norm class is the square class of
$B(\alpha_i,\alpha_i)$ in $K^\times/(K^\times)^2$; rescaling by
$t_i\in K^\times$ changes its norm by $t_i^2$ and preserves that class.

The golden-lattice applications also retain their proved inclusion
$\Gamma(P)\subseteq\operatorname{Ref}(L)$ for the literal invariant
lattice $L$ of Bugaenko's theorem~\cite{Bugaenko1984}. Here $\Gamma(P)$
is the original reflection group and $\operatorname{Ref}(L)$ the full
reflection group of that specified lattice. Equality of defining
fields alone would not give this inclusion.

Unlike the compact construction above, the arithmetic instance begins
with proved necessary covers at ranks $16$ and $17$. Their pruning is
an input; the LP supplies their final energy accounting.

\begin{assumption}[Arithmetic input package]
\label{a:input-package}
We use the following proved statements of Supplement~A on their complete
stated domains:
\begin{enumerate}[label=\textup{(\roman*)},leftmargin=*]
\item The actual common totally real field, original root lines,
 integral Cartan-cycle conditions, field-square restrictions, and
 positivity at every nonphysical real embedding are retained wherever
 they occur in a local theorem.
\item The complete necessary rank-seventeen and rank-sixteen covers
 are the fourteen and nineteen groups stated in
 Tables~\ref{a:rank17} and~\ref{a:rank16}, including all defect,
 spectator and field alternatives used to prove those covers.
\item The component and joint-component energy theorems give the rows
 in Section~\ref{a:costs} and Tables~\ref{a:rank17}--\ref{a:rank16}
 on the specified supports and complete forest domains. Each component
 floor already includes all contributions used in its local proof.
\item The complete local contradictions in Table~\ref{a:local-exits}
 hold on their indicated actual domains, including their boundary
 cases and the identity of each actual rank-one opposite.
\item The four computed outputs in
 Table~\ref{a:computed-inputs} hold at their proved complete input
 domains, with the original reduction and domain-inclusion theorems.
\item Every golden-lattice application retains the specified lattice
 and reflection-group inclusion above.
\end{enumerate}
\end{assumption}

This package specifies proved inputs of Supplement~A, rather than
consequences of real matrix positivity alone.

\begin{table}[htbp]
\centering
\caption{The four retained computed-output premises. A surviving model
is a necessary algebraic possibility, not an asserted polytope.}
\label{a:computed-inputs}
\begin{tabular}{>{\raggedright\arraybackslash}p{.08\linewidth}
 >{\raggedright\arraybackslash}p{.43\linewidth}
 >{\raggedright\arraybackslash}p{.39\linewidth}}
\toprule
Output&Complete reduced domain&Conclusion used by the local proof\\\midrule
(i)&The quadratic-field whole-$A_{17}$ domain&Empty\\
(ii)&The two cubic-field whole-$A_{17}$ finite-anchor domains&Empty
 in each of the two complete domains\\
(iii)&The shared quadratic pair domain reached separately from whole
 $B_{17}$ and whole $D_{17}$&Two necessary pairs, each retained until
 the final $B$- or $D$-specific energy estimate is applied\\
(iv)&The quadratic pair domain reached from whole $B_{16}$&One necessary
 pair, retained until its actual-edge energy estimate is applied\\\bottomrule
\end{tabular}
\end{table}

The whole-$D_{16}$ exclusion is a separate manual dyadic argument,
not a census premise. Its domain also differs from $D_{16}+I$ at
rank seventeen; both appear separately below.

\subsection{Notation and the arithmetic local rows}
\label{a:costs}

All component energies in this section are still projections of the
original vector $x$. Write $T_2$ for an ordinary irreducible $I_2(m)$,
$m\ge3$, $T_3$ for $A_3,B_3$ or $H_3$, and $T_4$ for one of
$A_4,B_4,D_4,F_4,H_4$. Let $\mathcal K$ denote either actual typed profile
$N_7$ or $O_7$, with the two choices retained as distinct branches.
An unqualified $A_1$ is ordinary and nonshort. The symbols $I,J,N_7,O_7$
retain the actual-wall definitions of Section~\ref{c:small}.

We use the following complete rank-five families:
\begin{align*}
\mathcal R_5={}&\{A_5,D_5,A_4+A_1,H_4+A_1,
                   T_3+T_2,2T_2+A_1\},\\
\mathcal R'_5={}&\{A_5,D_5,A_4+A_1,H_4+A_1,
                   A_3+T_2,2T_2+A_1\}.
\end{align*}
Every member has ordinary energy at least $5/14$ by
Table~\ref{c:ordinary-table}. The restriction to $A_3$ in
$\mathcal R'_5$ belongs to the rank-sixteen input cover.

In addition to the ordinary estimates in
Table~\ref{c:ordinary-table}, the following joint bounds are used:
\begin{align}
 e_{I+T_2}&>21/100,& e_{I+T_3}&>2/7,&
 e_{I+T_4}&>23/70,\label{a:I-joints}\\
 e_{I+A_1}&>5051/25000=.20204.\label{a:I-A1}
\end{align}
These are the actual-short-profile theorems of Supplement~A.
We also retain the propagated floors $e_{T_2}>1/5$ beside either
typed profile and $e_{A_1}>2/15$ beside $O_7$.

Tables~\ref{a:rank17} and~\ref{a:rank16} state the remaining
arithmetic floors directly on their domains. Each entry is a retained
local theorem of Supplement~A for the complete forest in the middle
column, with every displayed dihedral and field alternative retained.
Separate summands denote separate support rows; a joint support is
named explicitly. A floor on a named large component concerns that
component alone, even when its local proof uses the remaining walls.
Thus $16/19$ on $A_{14}$ beside ordinary $T_2$ and
$.91286445198$ beside a typed profile are distinct input statements.
The ordinary $T_2$ floor $3/19$ covers $m\ge4$ by its arithmetic
dihedral theorem and $m=3$ by $319/2000>3/19$.
For typed $B/D$ components the single floors $.91859$ and $.93859$
already include the end-wall correction $.0487$; it consumes no
additional energy variable. Only the rows needed by the displayed
certificate are specified here; additional valid rows in the retained
model may receive weight zero.

\subsection{The complete rank-seventeen cover and its rows}
\label{a:rank17-section}

Table~\ref{a:rank17} gives all fourteen input groups. Braces and
semicolons denote unions of complete alternatives. X1--X13 are the
local contradiction inputs in Table~\ref{a:local-exits}.

\begin{longtable}{>{\raggedright\arraybackslash}p{.06\linewidth}
 >{\raggedright\arraybackslash}p{.35\linewidth}
 >{\raggedright\arraybackslash}p{.49\linewidth}}
\caption{All rank-seventeen groups and their LP rows.}
\label{a:rank17}\\
\toprule
No.&Complete forest family&Rows giving total energy at least one,
 or retained local contradiction\\\midrule
\endfirsthead
\toprule No.&Complete forest family&Energy rows or contradiction\\\midrule\endhead
1&$A_{11}+I+(A_5\text{ or }D_5)$&
 $13/20+5/14=141/140$; the nonnegative $I$ energy is unused\\
2&$A_{11}+\mathcal K+U_4$, $U_4\in\{A_4,B_4,D_4,F_4\}$&
 $13/20+9/100+2/7$\\
3&$A_{12}+R$, $R\in\mathcal R_5$&
 $5/7+5/14$\\
4&$A_{12}$ plus $I+T_4$; $I+T_3+A_1$; or $I+2T_2$&
 $5/7$ plus, respectively, the joint $23/70$;
 the joint $2/7$ and $1/10$; or the joint $21/100$ and $1/7$\\
5&$A_{14}$ plus $T_3$; $T_2+A_1$; or $I+T_2$&
 $4/5$ plus, respectively, $1/4$; $1/7+1/10$;
 or the joint $21/100$\\
6&$A_{15}$ plus $T_2$; $2A_1$; $I+A_1$; $N_7$; or $O_7$&
 $T_2$: X1. The next two give $.8348+.2$ and
 $.8348+.20204$. Either typed branch gives $.9191408+.09$\\
7&$A_{16}$ plus ordinary $A_1$ or actual $I$&X2 in each
 of its complete rank-one alternatives\\
8&Entire $A_{17}$&X3, retaining outputs (i) and (ii) and their
 actual endpoint-domain inclusions\\
9&$B_{15}+N_7$ or $B_{15}+O_7$&
 $.93859+.09=1.02859$\\
10&$B_{16}+I$&X4\\
11&Entire $B_{17}$&$e_{B_{17}}\ge1.00184$, after the
 separate reduction and energy estimates for both survivors\\
12&$D_{15}+N_7$ or $D_{15}+O_7$&
 $.91859+.09=1.00859$\\
13&$D_{16}$ plus ordinary $A_1$ or actual $I$&
 X5 for ordinary $A_1$; X6 for actual $I$\\
14&Entire $D_{17}$&$e_{D_{17}}\ge1.00184$, after its
 genuine-fork reduction and its energy estimates for both survivors\\\bottomrule
\end{longtable}

The fourteen groups yield $39$ complete parameter-family branches,
without asserting their geometric realizability. Repeated symbols
such as $2T_2$ represent independently parameterized components.
In $I+2T_2$, choose one dihedral for the joint bound and give the
other its ordinary row; the $I$ energy is used once.

\subsection{The complete rank-sixteen cover and its rows}
\label{a:rank16-section}

Table~\ref{a:rank16} gives the separate nineteen-group rank-sixteen
input cover, including its narrower $\mathcal R'_5$ family.

\begin{longtable}{>{\raggedright\arraybackslash}p{.06\linewidth}
 >{\raggedright\arraybackslash}p{.35\linewidth}
 >{\raggedright\arraybackslash}p{.49\linewidth}}
\caption{All rank-sixteen groups and their LP rows.}
\label{a:rank16}\\
\toprule
No.&Complete forest family&Rows giving total energy at least one,
 or retained local contradiction\\\midrule
\endfirsthead
\toprule No.&Complete forest family&Energy rows or contradiction\\\midrule\endhead
1&$A_{10}+I+(A_5\text{ or }D_5)$&$404/625+5/14$\\
2&$A_{10}+\mathcal K+U_4$, $U_4\in\{A_4,B_4,D_4,F_4\}$&
 $404/625+9/100+2/7$\\
3&$A_{11}+R$, $R\in\mathcal R'_5$&$13/20+5/14=141/140$\\
4&$A_{11}+I+T_4$&The joint $A_{11}+I$ bound $18/25$,
 plus the disjoint $T_4$ bound $2/7$, giving $176/175$\\
5&$A_{11}+O_7+A_3$&$13/20+9/100+13/50=1$\\
6&$A_{12}$ plus $A_4$; $H_4$; or $2T_2$&$5/7+2/7=1$\\
7&$A_{12}+I+T_3$&$5/7$ plus the joint $I+T_3$ bound $2/7$\\
8&$A_{12}+\mathcal K+T_2$&$5/7+9/100+1/5$, with the actual
 typed-profile propagation on the last component\\
9&$A_{13}$ plus $A_3$; $T_2+A_1$; $I+T_2$; or $O_7+A_1$&
 Respectively: $53/70+1/4$; $53/70+1/7+1/10=1$;
 X7; or $233/300+9/100+2/15=1$\\
10&$A_{14}$ plus $T_2$; $2A_1$; $I+A_1$; $N_7$; or $O_7$&
 Respectively: $16/19+3/19=1$; $4/5+.2=1$;
 $4/5+.20204$; or $.91286445198+.09$ for either typed branch\\
11&$A_{15}$ plus ordinary $A_1$ or actual $I$&X8 in its
 complete isolated/nonisolated rank-one domains\\
12&Entire $A_{16}$&X9\\
13&$B_{14}+N_7$ or $B_{14}+O_7$&
 $.91859+.09=1.00859$\\
14&$B_{15}$ plus ordinary $A_1$ or actual $I$&
 X10 for ordinary $A_1$; X11 for actual $I$\\
15&Entire $B_{16}$&$e_{B_{16}}\ge1.00640742675$, after
 its own domain inclusion and edge estimate for the unique survivor\\
16&$D_{13}$ plus $I+T_2$ or $O_7+A_1$&
 $.82436$ plus the remainder floor $.21$; joint on $I+T_2$,
 and coarsened from $.09+2/15$ on $O_7+A_1$\\
17&$D_{14}$ plus ordinary $I_2(m)$, $m\ge4$; $N_7$; or $O_7$&
 Ordinary: $.84816+3/19$. Either typed branch:
 $.91859+.09$\\
18&$D_{15}$ plus ordinary $A_1$ or actual $I$&
 X10 for ordinary $A_1$; X12 for actual $I$\\
19&Entire $D_{16}$&X13, the actual quadratic reduction and
 complete dyadic argument\\\bottomrule
\end{longtable}

Each sum is at least one, which suffices against the strict budget
\eqref{c:budget}; no uniform positive excess is needed. For example,
group 5 uses three disjoint rows summing to exactly one. Group 4
instead uses the joint support $A_{11}+I$:
\[
 e_{A_{11}}+e_I\ge18/25,\qquad e_{T_4}\ge2/7.
\]
These rows have disjoint supports. The floor $18/25$ must not be
assigned to $A_{11}$ alone or combined with a second full use of $I$.

The nineteen groups give $49$ branches. Together the two covers have
$88$: $72$ energy branches and $16$ local contradiction branches,
with unbounded dihedral parameters quantified within their full domains.

\subsection{The arithmetic local contradiction statements}
\label{a:local-exit-section}

For completeness, Table~\ref{a:local-exits} specifies every local
contradiction used by Tables~\ref{a:rank17} and~\ref{a:rank16}. These
are imported local theorems with their full arithmetic and geometric
hypotheses. They are not extra scalar energy claims. Their LP row is
zero mass on the exactly matched domain.

\begin{longtable}{>{\raggedright\arraybackslash}p{.07\linewidth}
 >{\raggedright\arraybackslash}p{.31\linewidth}
 >{\raggedright\arraybackslash}p{.52\linewidth}}
\caption{Arithmetic local contradiction inputs. ``Entire'' refers to
the complete original vertex forest.}\label{a:local-exits}\\
\toprule
Code&Matched entire forest&Complete local argument retained\\\midrule
\endfirsthead
\toprule Code&Matched entire forest&Local argument\\\midrule\endhead
X1&$A_{15}+T_2$, rank $17$&The ordinary-dihedral theorem,
 including every silver, golden and low-energy label domain.\\
X2&$A_{16}+A_1$ or $A_{16}+I$, rank $17$&The complete
 actual rank-one-opposite theorem, isolated and nonisolated.\\
X3&Entire $A_{17}$, rank $17$&The endpoint and field reductions,
 with proved finite-anchor inclusions in outputs (i) and (ii).\\
X4&$B_{16}+I$, rank $17$&The complete six-role short-cap theorem,
 with its quadratic and cubic field alternatives.\\
X5&$D_{16}+A_1$, rank $17$, with ordinary $A_1$&The ordinary
 rank-one theorem, retaining both fork alternatives.\\
X6&$D_{16}+I$, rank $17$&The actual root-line theorem in degrees
 two, three and four, including both signed forks and all dyadic cases.\\
X7&$A_{13}+I+T_2$, rank $16$&The joint endpoint-and-remainder
 theorem with the complete ordinary dihedral domain.\\
X8&$A_{15}+A_1$ or $A_{15}+I$, rank $16$&The complete
 actual rank-one-opposite theorem, with all height boundaries.\\
X9&Entire $A_{16}$, rank $16$&The endpoint, field, finite-anchor
 and time-ratio theorem for this whole component.\\
X10&$B_{15}+A_1$ or $D_{15}+A_1$, rank $16$, with ordinary $A_1$&
 Their respective nonisolated reductions and ordinary isolated theorem,
 including $h=1/2$.\\
X11&$B_{15}+I$, rank $16$&The complete isolated-short theorem
 for $0<h<1/2$.\\
X12&$D_{15}+I$, rank $16$&The complete isolated-short theorem
 for $0<h<1/2$, including both signed forks.\\
X13&Entire $D_{16}$, rank $16$&The quadratic reduction and manual
 dyadic theorem: odd, inert, ramified and split cases.\\\bottomrule
\end{longtable}

The rank-one theorems follow the same actual opposite wall. Nonzero
contact with the large component gives the nonisolated alternative;
otherwise its actual isolated height lies in $0<h<1/2$ or $h\ge1/2$.
The equality case belongs to the ordinary theorem. These alternatives
cover an existing wall without replacing it by a formal extension.

The whole-$B_{17}$ and whole-$D_{17}$ energy rows use separate proved
inclusions in output (iii) and separate energy estimates for its two
survivors. Output (iv) likewise retains its unique pair until the actual
$B_{16}$ edge estimate is applied. These outputs are not emptiness
claims; whole $D_{16}$ instead uses the manual theorem X13.

\subsection{One joined arithmetic certificate}
\label{a:certificate}

The arithmetic tables give a hand certificate. On each energy branch in
Tables~\ref{a:rank17} and~\ref{a:rank16}, give the displayed disjoint
support rows weight one and all additional rows weight zero. Each
original energy has total weight at most one, and the weighted lower
bounds sum to at least one. A branch X1--X13 has zero mass and hence
zero normalized deficit by its retained local theorem.

The actual computation uses the same assembler for all $88$ arithmetic
branches and for the $829$ compact branches at dimensions $18$--$29$.
The joined catalogue has $917$ branches, including $167$ local
contradiction domains; its checked capacity dual has $3324$ variables
and $2488$ rows. The reproduction section records these global LP
checks separately from the retained local proofs.

\begin{theorem}[Compact arithmetic recovery]\label{a:recovery}
Retain Input~\ref{a:input-package}, the compact local inputs of
Section~\ref{c:section}, and the published compact ceiling $29$.
The joined energy LP is infeasible for every actual dimension at least
$16$. It recovers the existing compact arithmetic bound $n\le15$.
\end{theorem}
\begin{proof}
Every actual arithmetic configuration at ranks $16$ and $17$ belongs
to the complete corresponding cover. A branch with a retained local
contradiction is impossible. On every other branch the exact capacity
certificate gives total original energy at least one, contrary to
\eqref{c:budget}. The compact domains at ranks $18$--$29$ apply unchanged,
since an arithmetic compact polytope is compact; their certificate was
proved independently of arithmetic in Theorem~\ref{c:recovery}.
The external ceiling excludes dimensions $30$ and above.
\end{proof}

